% ---- pg-409 （书页 397） ----

\noindent\textbf{例 9.10}\quad (股票价格数据的旋转载荷)

\hyperref[tab:9.8]{表 9.8} 表示对例 8.5 和例 9.5 的股票价格数据的因子载荷的初始值估计和极大似然估计. 假定 $m=2$ 的因子模型, 特殊方差估计和由各因子解释的(标准化)样本总方差的累积比例, 也已给出.

由未旋转载荷所提出的因子阐述, 已在例 9.5 给出. 我们确认过市场因子和行业因子.

旋转载荷表明银行股票(J.P. Morgan, Citibank 和 Wells Fargo)在第一因子上高载荷, 而石油股票(Royal Dutch Shell 和 ExxonMobil)在第二因子上高载荷. (虽然没显示从主成分分解得到的旋转载荷, 对它们也能观测到相同现象.) 这两个旋转因子一起将行业区分开来. 想标明这些因子是困难的. 因子 1 代表那些引起银行股票一起波动的独特经济力量, 因子 2 代表影响石油类股票的经济条件.

如我们已经注意到的, 一个一般因子(也即所有变量在其上都有高载荷的因子)倾向于在旋转后被破坏掉了. 为了这个缘故, 在出现一般因子的情况下, 在进行正交旋转时, 有时将一般因子载荷固定住.\footnote{{\footnotesize 某些一般目的因子分析程序允许人们固定与某些因子有关的载荷而旋转其余因子.}} \hfill$\blacksquare$

\begin{mvtable}
\begin{tabular}{lccccc}
\toprule
\multirow{2}{*}{变量} & \multicolumn{2}{c}{\shortstack{因子载荷的\\极大似然估计}} & \multicolumn{2}{c}{因子载荷的旋转估计} & \multirow{2}{*}{特殊方差 $\hat{\psi}_i^2=1-\hat{h}_i^2$}\\
\cmidrule(lr){2-3}\cmidrule(lr){4-5}
& $F_1$ & $F_2$ & $F_1^*$ & $F_2^*$ & \\
\midrule
1.\ J.\ P.\ Morgan & 0.115 & 0.755 & \textbf{0.763} & 0.024 & 0.42\\
2.\ Citibank & 0.322 & 0.788 & 0.821 & 0.227 & 0.27\\
3.\ Wells Fargo & 0.182 & 0.652 & \textbf{0.669} & 0.104 & 0.54\\
4.\ Royal Dutch Shell & 1.000 & $-0.000$ & 0.118 & \textbf{0.993} & 0.00\\
5.\ ExxonMobil & 0.683 & 0.032 & 0.113 & \textbf{0.675} & 0.53\\
\midrule
\shortstack{所解释的\\样本总方差的\\累积比例} & 0.323 & 0.647 & 0.346 & 0.647 & \\
\bottomrule
\end{tabular}
\phantomsection\label{tab:9.8}\mvtabcap{表 9.8}
\end{mvtable}

\noindent\textbf{例 9.11}\quad (奥林匹克十项全能运动数据的旋转载荷)

例 9.6 中对奥林匹克十项全能运动数据给出了估计的因子载荷和特殊方差. 对一个 $m=4$ 的因子模型, 既用主成分法也用极大似然解法, 导出了这些量. 所有隐含因子的阐述并未马上就清楚明了. 为弄清旋转变量载荷是否会进一步加深认识, 我们做过最大方差旋转[见\hyperref[eq:9-45]{式(9-45)}]. \hyperref[tab:9.9]{表 9.9} 中除了特殊方差外, 还有对 $m=4$ 因子解的最大方差旋转载荷. 略去估计载荷, 旋转只影响由各因子解释的样本总方差之比的分布. 由所有因子解释的样本总方差之累积比, 没有变.

对两种解法的旋转因子载荷表明相同的潜在性质, 虽然因子 1 和 2 不在同一次序上. 我们看到掷铅球、掷铁饼和掷标枪在一个因子上有高载荷, 并且依照林登(文献[11])的意见称此因子为爆发性臂力强度. 类似地, 跳高、110 米栏、撑竿跳, 以及——在某种程度上——跳远, 都在另一因子上有高载荷. 林登把这个因子标为爆发性腿部强度. 100 米、400 米, 以及——还是在某种程度上——跳远在第三个因子上有高载荷, 这个因子可以叫做跑的速度因子. 最后,
